\chapter{主观题}
\begin{enumerate}
  \item 题型设计。
        \begin{enumerate}
          \item 案情设计。
                \begin{enumerate}
                  \item 从字数。较短（400 - 500）、中等（700 - 800）、\zx{超长}（1500，\al{19 年}）。
                  \item 内容连接性。各段独立、前后联系、
                \end{enumerate}
          \item 问题设计。
                \begin{enumerate}
                  \item \zx{一问一答}。3 - 6 个问题。\al{18 年}。
                  \item \zx{一句话问法}。\al{04 年}。
                \end{enumerate}
        \end{enumerate}
  \item 答题要求和写作模板。哪些内容才是需要回答的。
        \begin{enumerate}
          \item 内容要求。
                \begin{enumerate}
                  \item 犯罪形态。一般情形不要答，特殊情形（未完成形态）才需要答。
                        \begin{enumerate}
                          \item 就是在考既遂。
                          \item 共同犯罪。
                          \item 中途参与、事后参与。
                        \end{enumerate}
                  \item 共同犯罪。\bs{先实行犯，后教唆犯、帮助犯}。
                        \begin{enumerate}
                          \item 特殊要答，一般情形不答。\al{直接正犯 \cuo。间接正犯、共同正犯 \dui}。
                          \item 主要考定性、不考量刑。\al{主犯、从犯 \cuo。胁从犯 \dui}。
                        \end{enumerate}
                  \item 罪数。\al{想象竞合、牵连犯、吸收犯、不可罚的事后行为 \dui}。
                  \item 数额。问到才答。
                  \item 量刑。主要考自首、立功等量刑制度，不考具体建议。\al{应当从重、从轻 \cuo}。
                  \item 具体法条。不需要。
                  \item 罪名。认为应当构成 A 罪时，才考虑是否回答不构成 B 罪。\al{贷款诈骗、骗取贷款}。
                \end{enumerate}
          \item 形式要求。
                \begin{enumerate}
                  \item 一级标题、二级标题。
                        \begin{enumerate}
                          \item 案情长。事实 \yjt 人头 \yjt 共同犯罪的部分 \yjt 单独犯罪的部分。
                          \item 案情短。人头 \yjt 行为。
                        \end{enumerate}
                  \item \bs{先结论，后理由}。
                \end{enumerate}
        \end{enumerate}
  \item 考察能力。
        \begin{enumerate}
          \item 阅读能力。\zd{发现考点，提炼考点的能力}。
                \begin{enumerate}
                  \item 人物 + 线索。
                  \item 人物 + 行为。
                  \item 要求。不要漏点、读一遍能反映出啥考点、太长的读一段分析一段写一段。
                  \item \bs{理不清。找被害人}。
                \end{enumerate}
          \item 写作能力。文字组织、答题结构。
        \end{enumerate}
  \item 论证说理和观点展示。逻辑论证能力。
        \begin{enumerate}
          \item 分则。\al{是否构成 A 罪 \yjt A 罪的\xhx{关键}构成要件}。
          \item 总则。核心还是概念。\al{犯罪预备 \yjt 是否着手}。
          \item 关联说理。问题具有关联性。\al{拿卡、用卡}。
        \end{enumerate}
  \item \zx{观点展示}。\al{2021}。
        \begin{enumerate}
          \item 学术型（数量有限）、实务型（太多了，怎么应用标准）。\al{钱放到垃圾桶 \yjt 取得、控制}。
                \begin{enumerate}
                  \item 学术型汇总。
                        \begin{enumerate}
                          \item 作为义务的来源根据。考不足、缺点。
                          \item 因果关系判断方法。
                          \item 偶然防卫、偶然避险。\zd{故意型}。
                          \item 事后防卫。\al{龙哥案}。
                          \item 防卫过当。
                          \item 打击错误。
                          \item 着手。犯意表现（认定的太早）；实施构成要件；寄毒药，寄出着手，收到着手。
                          \item 不能犯、未遂犯。主观有就有；具体危险说、客观危险说（事后科学鉴定）。
                          \item 犯罪中止、未遂。
                          \item 共同犯罪。过失的共同正犯。\al{打野猪}。
                          \item 共犯从属性、独立性。
                          \item 片面帮助、实行。
                        \end{enumerate}
                  \item 实务型汇总。逻辑推导。
                        \begin{enumerate}
                          \item 被害人自陷风险。\al{放火进去救人、救钱；暴风雨摆渡}。
                          \item 作为义务。\al{防卫行为义务来源}。
                          \item 结果加重犯。
                          \item 防卫对象。
                          \item 避险过当。
                          \item 事前故意。
                          \item 结果的提前发生。
                          \item 重复侵害行为。
                          \item 中立的帮助行为。
                          \item 承继的共犯。
                          \item 连续犯。
                          \item 教唆他人自杀。
                          \item 婚内强奸。
                          \item 非法拘禁、故意伤害。
                          \item 绑架既遂。
                          \item 绑架 + 杀人。
                          \item 拐卖妇女。被害人承诺被卖。
                          \item 财产犯罪所有权。黑吃黑。
                          \item 不法原因给付物。
                          \item 死者占有。
                          \item 盗窃房产产权。产权归属，形式、实质。
                          \item 错误汇款。
                          \item 盗窃秘密性。
                          \item 事后转化抢劫。暴力对象、既遂标准。
                          \item 冒充军警抢劫。
                          \item 敲诈勒索和行使权利。
                        \end{enumerate}
                \end{enumerate}
          \item 观点推导过程。\al{贿赂款性质}。
          \item 答 2 种即可，正反。
        \end{enumerate}
  \item 常考知识点汇总。
\end{enumerate}
